\week{W9}{Control 2: computed torque, Cartesian, force and hybrid control}{9}
% ---- local macros (prefix wn) ----
\newif\ifwnAq
\newenvironment{wnfig}{\par\noindent\begin{minipage}{\linewidth}}{\end{minipage}\par}
\newcommand{\wnE}{E}\newcommand{\wnEd}{\dot E}
% generic partitioned-controller diagram (slides p.11, p.20); set \wn* then call \wnloop
\newcommand{\wnloop}{\begin{tikzpicture}[x=0.85cm,y=0.68cm]
\node[sgn,anchor=east] (i2) at (0,0) {$\wnDD$};
\node[sgn,anchor=east] (i1) at (0,-2.5) {$\wnD$};
\node[sgn,anchor=east] (i0) at (0,-3.5) {$\wnR$};
\node[sumj] (ss) at (2.1,0) {};
\node[svblk,minimum width=9mm] (kp) at (1.2,-1.3) {$\wnKp$};
\node[svblk,minimum width=9mm] (kv) at (2.8,-1.3) {$\wnKv$};
\node[sumj] (sv) at (2.8,-2.5) {};
\node[sumj] (sp) at (1.2,-3.5) {};
\node[mbblk] (al) at (4.7,0) {$\wnAl$};
\node[sumj] (sb) at (5.9,0) {};
\node[mbblk] (be) at (7.95,-1.3) {$\wnBe$};
\node[plant] (rb) at (8.1,0) {$\wnPl$};
\draw[sig] (i2.east) -- (ss.west); \node[sgn] at (1.6,0.22) {$+$};
\draw[sig] (kp.north) -- (ss.235); \node[sgn] at (1.35,-0.55) {$+$};
\draw[sig] (kv.north) -- (ss.305); \node[sgn] at (2.65,-0.55) {$+$};
\draw[sig] (ss) -- node[sgn,above]{$\wnFp$} (al);
\draw[sig] (al) -- node[sgn,pos=0.45,above]{$+$} (sb);
\draw[sig] (be.west) -- (5.9,-1.3) -- (sb.south); \node[sgn] at (6.12,-0.55) {$+$};
\draw[sig] (sb) -- node[sgn,above]{$\wnU$} (rb);
\draw[sig] (rb.east) -- (11.0,0) node[sgn,right]{$\wnP,\wnV$};
\draw[line width=0.6pt] (10.2,0) -- (10.2,-3.5);
\fill (10.2,0) circle (0.9pt); \fill (10.2,-1.3) circle (0.9pt); \fill (10.2,-2.5) circle (0.9pt);
\draw[sig] (10.2,-1.3) -- (be.east);
\draw[sig] (10.2,-2.5) -- node[sgn,above,pos=0.6]{$\wnV$} node[sgn,pos=0.93,above]{$-$} (sv.east);
\draw[sig] (10.2,-3.5) -- node[sgn,above,pos=0.6]{$\wnP$} node[sgn,pos=0.93,above]{$-$} (sp.east);
\draw[sig] (i1.east) -- (1.05,-2.5) arc[start angle=180,end angle=0,radius=0.15] -- (sv.west); \node[sgn] at (2.35,-2.28) {$+$};
\draw[sig] (i0.east) -- (sp.west); \node[sgn] at (0.8,-3.28) {$+$};
\draw[sig] (sv.north) -- node[sgn,right]{$\wnEd$} (kv.south);
\draw[sig] (sp.north) -- node[sgn,left]{$\wnE$} (kp.south);
\ifwnAq\draw[sig] (10.2,0) -- (10.2,1.0) -- (4.7,1.0) -- (al.north); \node[sgn,above] at (7.4,1.0) {$\wnP$};\fi
\node[grpsv,fit=(ss)(kp)(kv)(sv)(sp)] (gsv) {};
\node[grplab,text=svcol,anchor=south west] at (gsv.north west) {SERVO};
\node[grpmb,fit=(al)(sb)(be)(rb)] (gmb) {};
\node[grplab,text=mbcol,anchor=north east,yshift=2pt] at (gmb.south east) {MODEL-BASED};
\end{tikzpicture}}

% ============================================================
\begin{multicols}{3}
\hd{Which law for which problem? \added}
{\footnotesize
\begin{tabularx}{\linewidth}{@{}>{\raggedright\arraybackslash}p{0.21\linewidth}>{\raggedright\arraybackslash}p{0.26\linewidth}X@{}}
\toprule
Law & needs & gives / use\\\midrule
joint PD/PID (W8) & none or $\mathbf g$; $\mathbf q,\dot{\mathbf q}$ & stable; slow/geared; steady error without integral\\
computed torque \pg{18--21} & exact $\mathbf M,\mathbf C,\mathbf g,\mathbf f$ & $\ddot{\mathbf E}+\mathbf K_v\dot{\mathbf E}+\mathbf K_p\mathbf E=\mathbf 0$, decoupled\\
inexact model \pg{27--28} & $\hat{\mathbf M},\hat{\mathbf C},\hat{\mathbf g},\hat{\mathbf f}$ & forced error eq.: small, not zero\\
Cartesian \pg{32--33} & $\mathbf M_x,\mathbf c_x,\mathbf g_x,\mathcal K,\mathbf J$ & same form in $\mathbf x$; no IK or $\mathbf J^{-1}$; fails at singularity\\
force \pg{36--43} & $m$ ($k_e$); $f_e,\dot x$ & Opt.\,2 $E_{ss}=f_{dist}/(1+mk_e^{-1}k_p)$\\
hybrid \pg{47} & Cartesian + force law & each axis obeys one law\\
\bottomrule
\end{tabularx}}\par
\textbf{Go to:} write law p.\,\pageref{w9:rec}; $K_p,K_v$ p.\,\pageref{w9:gain}; ``mass 10\,\% off'' p.\,\pageref{w9:ten}; Cartesian p.\,\pageref{w9:cart}; force/wall p.\,\pageref{w9:force}; which axes p.\,\pageref{w9:hyb}; practical sheet p.\,\pageref{w9:prac}.

\hd{Scalar partitioned law: cubic spring, single link \pg{7--15}}
\premise{$m\ddot x+b\dot x+kx^3=f$; $x_d,\dot x_d,\ddot x_d$; measured $x,\dot x$}{$f$ giving $\ddot E+k_v\dot E+k_pE=0$}
\begin{fbx}[Partitioned (computed-torque) law, scalar \pg{10}]
\[f=\alpha\,\acute f+\beta,\quad \alpha=m,\ \ \beta=b\dot x+kx^3,\quad \acute f=\ddot x_d+k_v\dot E+k_pE,\quad E=x_d-x\]
where $\alpha$ = coefficient of $\ddot x$ (kg); $\beta$ = the rest, with the \emph{measured} $x,\dot x$ (N); $\acute f$ = commanded acceleration (m/s$^2$); $k_p,k_v$ in s$^{-2}$, s$^{-1}$.
\end{fbx}
\begin{deriv}[closed loop, every line \pg{10}]
Insert $f$: $m\ddot x+b\dot x+kx^3=m\acute f+b\dot x+kx^3\Rightarrow\ddot x=\acute f$. With $\ddot E=\ddot x_d-\ddot x$:
\[\boxed{\ddot E+k_v\dot E+k_pE=0}\]
linear whatever the non-linearity; stable for $k_p,k_v>0$.
\end{deriv}
\begin{recipe}[Write a partitioned law \pg{15}]\label{w9:rec}
\begin{enumerate}
\item EoM as $(\text{coeff})\ddot q+(\text{rest})=\text{input}$ (Lagrange).
\item $\alpha$ = coeff of $\ddot q$; $\beta$ = rest, signs kept, measured state.
\item $\acute q=\ddot q_d+k_v\dot E+k_pE$, $E=q_d-q$.
\item input $=\alpha\acute q+\beta$; closed loop $\ddot E+k_v\dot E+k_pE=0$.
\item Check units; at rest on the path law $=\beta$ (e.g.\ $mgL\cos q$).
\end{enumerate}
\end{recipe}
\begin{worked}[cubic spring; $m=2$, $b=4$, $k=50$, $k_p=100$, $k_v=20$]
$x=0.10$, $\dot x=0.30$; $x_d=0.12$, $\dot x_d=0.40$, $\ddot x_d=0.50$. $E=0.02$, $\dot E=0.10$. $\acute f=0.50+20(0.10)+100(0.02)=4.5$; $\beta=4(0.30)+50(0.10)^3=1.25$ N; $f=2(4.5)+1.25=10.25$ N. Check $\ddot x=(10.25-1.25)/2=4.5=\acute f$ \checkmark. \chk
\end{worked}
\begin{trap}
$f=\acute f+\beta$ (no $\alpha$) mixes m/s$^2$ and N; $\ddot x_d$ goes in $\acute f$, not $\beta$; $E=x_d-x$.
\end{trap}
\begin{wnfig}
\wnAqfalse
\def\wnDD{\ddot x_d}\def\wnD{\dot x_d}\def\wnR{x_d}\def\wnKp{k_p}\def\wnKv{k_v}\def\wnAl{m}\def\wnBe{b\dot x+kx^3}\def\wnPl{m\ddot x{+}b\dot x{+}kx^3{=}f}\def\wnU{f}\def\wnP{x}\def\wnV{\dot x}\def\wnE{E}\def\wnEd{\dot E}\def\wnFp{\acute f}
\resizebox{\linewidth}{!}{\wnloop}\par
\drawncap{non-linear SISO partitioned controller (slide p.\,11). \textbf{Orange} = model-based ($m$, $\beta$), \textbf{blue} = servo ($k_p,k_v$ on the error), grey = plant. $\ddot x_d$ feed-forward; $\beta$ from measured $x,\dot x$.}
\end{wnfig}

\sh{Single link, viscous + Coulomb friction \pg{12--14}}
\premise{link $L$, tip mass $m$, angle $q$ from horizontal, friction $b_v,b_c$; $q_d,\dot q_d,\ddot q_d$}{torque $\tau$}
\begin{deriv}[EoM, slide p.\,13]
$\tau_v=b_v\dot q$, $\tau_c=b_c\operatorname{sgn}\dot q$. $x=L\cos q$, $y=L\sin q$: $K=\tfrac12mL^2\dot q^2$, $P=mgL\sin q$. $Q=\frac d{dt}\frac{\partial\mathcal L}{\partial\dot q}-\frac{\partial\mathcal L}{\partial q}$ with $\frac{\partial\mathcal L}{\partial\dot q}=mL^2\dot q$, $\frac{\partial\mathcal L}{\partial q}=-mgL\cos q$:
\[\boxed{\tau=mL^2\ddot q+b_v\dot q+b_c\operatorname{sgn}(\dot q)+mgL\cos q}\]
\end{deriv}
\begin{fbx}[Control law of the single link \pg{14}]
\[\tau=\alpha\acute\tau+\beta,\ \ \alpha=mL^2,\ \ \beta=b_v\dot q+b_c\operatorname{sgn}(\dot q)+mgL\cos q,\ \ \acute\tau=\ddot q_d+k_v\dot E+k_pE\]
$\alpha$ = inertia about the pivot (kg\,m$^2$); at rest on the path $\tau=mgL\cos q$.
\end{fbx}
\begin{worked}[$m=2$, $L=0.5$, $b_v=0.3$, $b_c=0.2$, $k_p=64$, $k_v=16$]
$q=30^\circ$, $\dot q=1$, $\ddot q_d=2$, $E=0.05$, $\dot E=-0.1$. $\alpha=0.5$; $\acute\tau=2+16(-0.1)+64(0.05)=3.6$; $\beta=0.3+0.2+2(9.81)(0.5)\cos30^\circ=8.9957$; $\tau=0.5(3.6)+8.9957=10.7957$ N\,m. Check $\ddot q=(\tau-\beta)/\alpha=3.6=\acute\tau$ \checkmark. \chk
\end{worked}
\begin{trap}
Degrees in $\cos q$; $\operatorname{sgn}$ of $\dot q$ (not $q$, $E$); $\alpha=mL$; $\cos q$ dropped. Real laws smooth $\operatorname{sgn}$ \added. Limit \pg{15}: needs the structure and parameters of the model.
\end{trap}

\hd{MIMO computed torque: $n$ joints \pg{16--21}}
\premise{$\mathbf Q=\mathbf M(\mathbf q)\ddot{\mathbf q}+\mathbf C(\mathbf q,\dot{\mathbf q})+\mathbf g(\mathbf q)+\mathbf f(\dot{\mathbf q})$; $\mathbf q_d,\dot{\mathbf q}_d,\ddot{\mathbf q}_d$; measured $\mathbf q,\dot{\mathbf q}$}{$\mathbf Q$ ($n\times1$)}
\begin{fbx}[Computed-torque law of a manipulator \pg{18--19}]
\[\begin{gathered}\mathbf Q=\boldsymbol\alpha\acute{\mathbf Q}+\boldsymbol\beta,\quad \boldsymbol\alpha=\mathbf M(\mathbf q),\quad \boldsymbol\beta=\mathbf C(\mathbf q,\dot{\mathbf q})+\mathbf g(\mathbf q)+\mathbf f(\dot{\mathbf q}),\\ \acute{\mathbf Q}=\ddot{\mathbf q}_d+\mathbf K_v\dot{\mathbf E}+\mathbf K_p\mathbf E,\quad \mathbf E=\mathbf q_d-\mathbf q\end{gathered}\]
where $\mathbf M$ $n\times n$ symmetric positive definite (kg\,m$^2$); $\mathbf C$ Coriolis/centripetal, $\mathbf g$ gravity, $\mathbf f$ friction (N\,m); $\mathbf K_p,\mathbf K_v$ \emph{diagonal} (s$^{-2}$, s$^{-1}$). $\boldsymbol\alpha=\mathbf M$ is \textbf{not} diagonal: it decouples.
\end{fbx}
\begin{deriv}[closed loop \pg{21}]
Insert $\mathbf Q$ in the plant: $\mathbf M\ddot{\mathbf q}+\mathbf C+\mathbf g+\mathbf f=\mathbf M\acute{\mathbf Q}+\mathbf C+\mathbf g+\mathbf f\Rightarrow\mathbf M\ddot{\mathbf q}=\mathbf M\acute{\mathbf Q}\Rightarrow\ddot{\mathbf q}=\acute{\mathbf Q}$ ($\mathbf M$ invertible). With $\ddot{\mathbf q}=\ddot{\mathbf q}_d-\ddot{\mathbf E}$:
\[\boxed{\ddot{\mathbf E}+\mathbf K_v\dot{\mathbf E}+\mathbf K_p\mathbf E=\mathbf 0}\ \Rightarrow\ \ddot e_i+k_{vi}\dot e_i+k_{pi}e_i=0\ \ (i=1..n)\]
$n$ decoupled linear equations; the off-diagonals of $\mathbf M$ cancel the coupling $M_{ij}\ddot q_j$.
\end{deriv}
\begin{wnfig}
\wnAqtrue
\def\wnDD{\ddot{\mathbf q}_d}\def\wnD{\dot{\mathbf q}_d}\def\wnR{\mathbf q_d}\def\wnKp{\mathbf K_p}\def\wnKv{\mathbf K_v}\def\wnAl{\mathbf M(\mathbf q)}\def\wnBe{\mathbf C(\mathbf q,\dot{\mathbf q}){+}\mathbf g(\mathbf q){+}\mathbf f(\dot{\mathbf q})}\def\wnPl{\text{Manipulator}}\def\wnU{\mathbf Q}\def\wnP{\mathbf q}\def\wnV{\dot{\mathbf q}}\def\wnE{\mathbf E}\def\wnEd{\dot{\mathbf E}}\def\wnFp{\acute{\mathbf Q}}
\resizebox{\linewidth}{!}{\wnloop}\par
\drawncap{MIMO computed torque (slide p.\,20). Slide's red box ``linearises'' = orange group; green ``stabilises'' = blue group. Arrow $\mathbf q\to\mathbf M(\mathbf q)$ added.}
\end{wnfig}
\begin{worked}[$\boldsymbol\alpha$ not diagonal: practical 3-link arm, $\mathbf q_0=(0,0,90^\circ)$]
$\mathbf M=\begin{bmatrix}103&31&1\\31&13&1\\1&1&1\end{bmatrix}$ (p.\,\pageref{w9:prac}). $\acute{\mathbf Q}=(1,0,0)$ needs $\boldsymbol\alpha\acute{\mathbf Q}=(103,31,1)$ N\,m: the $31,1$ hold joints 2, 3 still. Diagonal $\boldsymbol\alpha$ gives only $103$. \chk
\end{worked}
\begin{recipe}[Choose $\mathbf K_p,\mathbf K_v$, critical damping \added]\label{w9:gain}
\begin{enumerate}
\item Per joint $s^2+k_vs+k_p$; double root $-\omega_n$: $k_p=\omega_n^2$, $k_v=2\omega_n$ ($k_v^2=4k_p$).
\item $\omega_n$ from settling: $e=e_0(1+\omega_nt)e^{-\omega_nt}$ hits 2\,\% at $\omega_nt=5.83$, 5\,\% at $4.74$. \chk
\item $\mathbf K_p=\operatorname{diag}(\omega_i^2)$, $\mathbf K_v=\operatorname{diag}(2\omega_i)$ (act on accelerations, independent of inertia).
\item Check: equal negative roots; no overshoot.
\end{enumerate}
\end{recipe}
\begin{worked}[gains from settling time]
2\,\% in $0.5$ s: $\omega_n\ge5.83/0.5=11.67$; take $12$: $k_p=144$, $k_v=24$ ($24^2=4\cdot144$ \checkmark). \chk\ $\omega=(10,15,20)$: $\mathbf K_p=\operatorname{diag}(100,225,400)$, $\mathbf K_v=\operatorname{diag}(20,30,40)$. \chk\ $k_p=100,k_v=20$, $e(0)=0.1$: $e(0.1\,\text{s})=0.0736$. \chk
\end{worked}
\begin{trap}
``$\mathbf K$ diagonal $\Rightarrow\boldsymbol\alpha$ diagonal'': false. Error equation holds only for an exact model and a continuous law. Outer-loop gains (p.\,25) are torque gains.
\end{trap}

\hd{Imprecise model, discretisation, outer loop \pg{22--28}}
Model may be wrong (friction, wear, grasped load); the law runs in discrete time.
\begin{fbx}[Law and closed loop with hats \pg{27--28}]
\[\mathbf Q=\hat{\mathbf M}\acute{\mathbf Q}+\hat{\mathbf C}(\mathbf q,\dot{\mathbf q})+\hat{\mathbf g}(\mathbf q)+\hat{\mathbf f}(\mathbf q,\dot{\mathbf q}),\quad \acute{\mathbf Q}=\ddot{\mathbf q}_d+\mathbf K_v\dot{\mathbf E}+\mathbf K_p\mathbf E\]
\[\ddot{\mathbf E}+\mathbf K_v\dot{\mathbf E}+\mathbf K_p\mathbf E=\hat{\mathbf M}^{-1}\Big[(\mathbf M-\hat{\mathbf M})\ddot{\mathbf q}+(\mathbf C-\hat{\mathbf C})+(\mathbf g-\hat{\mathbf g})+(\mathbf f-\hat{\mathbf f})\Big]\]
Hat = controller's estimate. The right side is a \emph{forcing}: error does not decay to zero unless every hatted term is exact.
\end{fbx}
\begin{deriv}[the slide prints only the result \added]
Plant $=$ law: $\hat{\mathbf M}\acute{\mathbf Q}=\mathbf M\ddot{\mathbf q}+\boldsymbol\Delta$, $\boldsymbol\Delta=(\mathbf C-\hat{\mathbf C})+(\mathbf g-\hat{\mathbf g})+(\mathbf f-\hat{\mathbf f})$. So $\acute{\mathbf Q}-\ddot{\mathbf q}=\hat{\mathbf M}^{-1}(\mathbf M-\hat{\mathbf M})\ddot{\mathbf q}+\hat{\mathbf M}^{-1}\boldsymbol\Delta$, and $\acute{\mathbf Q}-\ddot{\mathbf q}=\ddot{\mathbf E}+\mathbf K_v\dot{\mathbf E}+\mathbf K_p\mathbf E$. \chk
\end{deriv}
\begin{fbx}[Single link, only the mass wrong \added]
$\hat\alpha=\hat mL^2$, $\hat\beta=\hat mgL\cos q$:
\[\ddot e+k_v\dot e+k_pe=\frac{m-\hat m}{\hat m}\Big[\ddot q+\frac gL\cos q\Big],\qquad e_{ss}\approx\frac{m-\hat m}{\hat m}\frac gL\frac{\cos q}{k_p}\]
$\hat m<m$: link sags, $e>0$; $e_{ss}=0$ at $90^\circ$. \chk
\end{fbx}
\begin{worked}[scenario: mass 10\,\% off; $m=2$, $L=0.5$, $k_p=100$, $k_v=20$, $g/L=19.62$]\label{w9:ten}
$\hat m=1.8$: ratio $0.2/1.8=1/9$. Held horizontal: forcing $\tfrac19(19.62)=2.18$; $e_{ss}=0.0218$ rad $=1.25^\circ$. At $30^\circ$: $0.0189$ rad (simulation $0.0191$). At $90^\circ$: $0$. $\hat m=2.2$: ratio $-1/11$, $e_{ss}=-0.0178$ rad $=-1.02^\circ$ (link rises). \chk\ Error $\approx$ forcing$/k_p$: not zero; $k_p\uparrow$ shrinks it but sampling limits $k_p$; integral action (W8) removes the constant part \added.
\end{worked}
\begin{slideerr}[p.\,19 vs 27]
Friction is $\mathbf f(\dot{\mathbf q})$ on p.\,19, $\hat{\mathbf f}(\mathbf q,\dot{\mathbf q})$ on pp.\,27--28: treat as $\mathbf f(\mathbf q,\dot{\mathbf q})$. Notation only.
\end{slideerr}
\textbf{Discretisation \pg{23--25}.} Computing $\mathbf M,\mathbf C,\mathbf g$ is slow, the PD part fast. \textbf{Fix:} evaluate the model on the \emph{planned} $\mathbf q_d,\dot{\mathbf q}_d,\ddot{\mathbf q}_d$ at a slow rate (even tabulated); the fast loop adds only $\mathbf K_p\mathbf E+\mathbf K_v\dot{\mathbf E}$.
\begin{wnfig}
\resizebox{\linewidth}{!}{\begin{tikzpicture}[x=0.85cm,y=0.68cm]
\node[sgn,anchor=east] (a2) at (0,1.5) {$\ddot{\mathbf q}_d$};
\node[sgn,anchor=east] (a1) at (0,0.9) {$\dot{\mathbf q}_d$};
\node[sgn,anchor=east] (a0) at (0,0.3) {$\mathbf q_d$};
\node[mbblk,minimum height=17mm,minimum width=22mm] (dm) at (2.6,0.9) {dynamics\\model};
\node[sumj] (sq) at (5.9,0.9) {};
\node[svblk,minimum width=9mm] (kp) at (5.1,-0.5) {$\mathbf K_p$};
\node[svblk,minimum width=9mm] (kv) at (6.7,-0.5) {$\mathbf K_v$};
\node[sumj] (sv) at (6.7,-1.7) {};
\node[sumj] (sp) at (5.1,-2.7) {};
\node[plant,minimum width=18mm] (rb) at (9.1,0.9) {manipulator};
\draw[sig] (a2.east) -- (dm.west |- a2);
\draw[sig] (a1.east) -- (dm.west |- a1);
\draw[sig] (a0.east) -- (dm.west |- a0);
\draw[sig] (dm.east) -- node[sgn,above,pos=0.45]{$\mathbf Q_{ff}$} (sq.west); \node[sgn] at (5.45,1.12) {$+$};
\draw[sig] (kp.north) -- (sq.235); \node[sgn] at (5.3,0.2) {$+$};
\draw[sig] (kv.north) -- (sq.305); \node[sgn] at (6.5,0.2) {$+$};
\draw[sig] (sq) -- node[sgn,above]{$\mathbf Q$} (rb);
\draw[sig] (rb.east) -- (11.2,0.9) node[sgn,right]{$\mathbf q,\dot{\mathbf q}$};
\draw[line width=0.6pt] (10.5,0.9) -- (10.5,-2.7); \fill (10.5,0.9) circle (0.9pt); \fill (10.5,-1.7) circle (0.9pt);
\draw[sig] (10.5,-1.7) -- node[sgn,above,pos=0.2]{$\dot{\mathbf q}$} node[sgn,above,pos=0.92]{$-$} (sv.east);
\draw[sig] (10.5,-2.7) -- node[sgn,above,pos=0.2]{$\mathbf q$} node[sgn,above,pos=0.92]{$-$} (sp.east);
\fill (1.0,0.9) circle (0.9pt); \fill (0.4,0.3) circle (0.9pt);
\draw[sig] (1.0,0.9) -- (1.0,0.45) arc[start angle=90,end angle=-90,radius=0.15] -- (1.0,-1.7) -- (4.95,-1.7) arc[start angle=180,end angle=0,radius=0.15] -- (sv.west); \node[sgn] at (6.35,-1.5) {$+$};
\draw[sig] (0.4,0.3) -- (0.4,-2.7) -- (sp.west); \node[sgn] at (4.6,-2.5) {$+$};
\draw[sig] (sv.north) -- node[sgn,right]{$\dot{\mathbf E}$} (kv.south);
\draw[sig] (sp.north) -- node[sgn,left]{$\mathbf E$} (kp.south);
\node[grpmb,fit=(dm)] (g1) {};
\node[grplab,text=mbcol,anchor=north west] at (g1.south west) {MODEL (slow)};
\node[grpsv,fit=(sq)(kp)(kv)(sv)(sp)] (g2) {};
\node[grplab,text=svcol,anchor=south west] at (g2.north west) {SERVO (fast)};
\end{tikzpicture}}\par
\drawncap{model in the outer loop (slide p.\,25). $\mathbf Q_{ff}$ is my name for the model output; the model reads only the plan.}
\end{wnfig}
\begin{slideerr}[p.\,25: gain units]
Servo terms are added \emph{to the torque}: $\mathbf K_p,\mathbf K_v$ are torque gains (N\,m/rad, N\,m\,s/rad), not s$^{-2}$, s$^{-1}$; only approximately decoupled \added. Practical link 1 ($\alpha=27$): $27\cdot100=2700$, $27\cdot20=540$. \chk
\end{slideerr}
Numbers: link 1 ($\alpha=27$, $\beta=88.29\cos q$, $k_p=100$, $k_v=20$), cubic $0\to\pi/2$ in $2$ s: continuous law $|e|<10^{-10}$; refresh every $10$ ms peak error $0.26$ mrad; $50$ ms $1.13$ mrad. \chk

\hd{Cartesian-space control \pg{29--33}}
\premise{tool path $\mathbf x_d,\dot{\mathbf x}_d,\ddot{\mathbf x}_d$ in $\{0\}$; FK $\mathcal K$, $\mathbf J(\mathbf q)$, dynamics}{$\mathbf Q$ so the \emph{tool} follows the path}
\textbf{Joint-space route \pg{30--31}:} $\mathbf q=\mathcal K^{-1}(\mathbf x)$, $\dot{\mathbf q}=\mathbf J^{-1}(\mathbf q)\dot{\mathbf x}$, $\ddot{\mathbf q}=\tfrac{d}{dt}\big(\mathbf J^{-1}(\mathbf q)\big)\dot{\mathbf x}+\mathbf J^{-1}(\mathbf q)\ddot{\mathbf x}$: IK and $\mathbf J^{-1}$ at every point, fails at a singularity.
\begin{slideerr}[p.\,31]
Slide prints $\dot J^{-1}\dot{\mathbf x}$. Correct: $\tfrac d{dt}\mathbf J^{-1}=-\mathbf J^{-1}\dot{\mathbf J}\mathbf J^{-1}$, so $\ddot{\mathbf q}=\mathbf J^{-1}(\ddot{\mathbf x}-\dot{\mathbf J}\dot{\mathbf q})$. \chk
\end{slideerr}
\begin{fbx}[Cartesian computed-torque law \pg{32--33}]\label{w9:cart}
\[\mathcal F=\mathbf M_x(\mathbf q)\ddot{\mathbf x}+\mathbf c_x(\mathbf q,\dot{\mathbf q})+\mathbf g_x(\mathbf q)\ \text{(plant)},\qquad \mathbf Q=\mathbf J^\top(\mathbf q)\mathcal F\]
\[\mathcal F=\mathbf M_x\big(\ddot{\mathbf x}_d+\mathbf K_v\dot{\mathbf E}_x+\mathbf K_p\mathbf E_x\big)+\mathbf c_x+\mathbf g_x,\quad \mathbf E_x=\mathbf x_d-\mathbf x\]
where $\mathcal F$ = tool force/moment in $\{0\}$ (N, N\,m); $\mathbf x=\mathcal K(\mathbf q)$, $\dot{\mathbf x}=\mathbf J\dot{\mathbf q}$; $\mathbf M_x$ tool inertia (kg), $\mathbf c_x,\mathbf g_x$ (N); $\mathbf K_p,\mathbf K_v$ (s$^{-2}$, s$^{-1}$). Closed loop $\ddot{\mathbf E}_x+\mathbf K_v\dot{\mathbf E}_x+\mathbf K_p\mathbf E_x=\mathbf 0$; no IK, no $\mathbf J^{-1}$ on the path; $\mathbf M_x$ blows up near a singularity.
\end{fbx}
\begin{deriv}[$\mathbf Q=\mathbf J^\top\mathcal F$ and $\mathbf M_x,\mathbf c_x,\mathbf g_x$ \added]
Virtual work: $\delta\mathbf x=\mathbf J\delta\mathbf q$, $\mathbf Q^\top\delta\mathbf q=\mathcal F^\top\mathbf J\delta\mathbf q\ \forall\delta\mathbf q\Rightarrow\mathbf Q=\mathbf J^\top\mathcal F$. Put into $\mathbf M\ddot{\mathbf q}+\mathbf C+\mathbf g=\mathbf Q$ with $\ddot{\mathbf q}=\mathbf J^{-1}(\ddot{\mathbf x}-\dot{\mathbf J}\dot{\mathbf q})$ (no friction):
\[\begin{gathered}\mathcal F=\underbrace{\mathbf J^{-\top}\mathbf M\mathbf J^{-1}}_{\mathbf M_x}\ddot{\mathbf x}+\underbrace{\mathbf J^{-\top}\big(\mathbf C-\mathbf M\mathbf J^{-1}\dot{\mathbf J}\dot{\mathbf q}\big)}_{\mathbf c_x}\\+\underbrace{\mathbf J^{-\top}\mathbf g}_{\mathbf g_x}\quad(\mathbf J\ \text{square, invertible})\end{gathered}\ \chk\]
\end{deriv}
\begin{wnfig}
\resizebox{\linewidth}{!}{\begin{tikzpicture}[x=0.84cm,y=0.68cm]
\node[sgn,anchor=east] (i2) at (0,0) {$\ddot{\mathbf x}_d$};
\node[sgn,anchor=east] (i1) at (0,-2.5) {$\dot{\mathbf x}_d$};
\node[sgn,anchor=east] (i0) at (0,-3.5) {$\mathbf x_d$};
\node[sumj] (ss) at (2.1,0) {};
\node[svblk,minimum width=9mm] (kp) at (1.2,-1.3) {$\mathbf K_p$};
\node[svblk,minimum width=9mm] (kv) at (2.8,-1.3) {$\mathbf K_v$};
\node[sumj] (sv) at (2.8,-2.5) {};
\node[sumj] (sp) at (1.2,-3.5) {};
\node[mbblk] (mx) at (4.7,0) {$\mathbf M_x(\mathbf q)$};
\node[sumj] (sb) at (5.9,0) {};
\node[mbblk,minimum width=11mm] (jt) at (7.3,0) {$\mathbf J^\top(\mathbf q)$};
\node[plant,minimum width=16mm] (rb) at (9.3,0) {manipulator};
\node[mbblk] (be) at (7.7,-1.3) {$\mathbf c_x(\mathbf q,\dot{\mathbf q}){+}\mathbf g_x(\mathbf q)$};
\node[mbblk,minimum width=10mm] (jf) at (9.6,-2.5) {$\mathbf J(\mathbf q)$};
\node[mbblk,minimum width=10mm] (kk) at (9.6,-3.5) {$\mathcal K(\ )$};
\draw[sig] (i2.east) -- (ss.west); \node[sgn] at (1.6,0.22) {$+$};
\draw[sig] (kp.north) -- (ss.235); \node[sgn] at (1.35,-0.55) {$+$};
\draw[sig] (kv.north) -- (ss.305); \node[sgn] at (2.65,-0.55) {$+$};
\draw[sig] (ss) -- (mx);
\draw[sig] (mx) -- node[sgn,pos=0.45,above]{$+$} (sb);
\draw[sig] (be.west) -- (5.9,-1.3) -- (sb.south); \node[sgn] at (6.12,-0.55) {$+$};
\draw[sig] (sb) -- node[sgn,above]{$\mathcal F$} (jt);
\draw[sig] (jt) -- node[sgn,above]{$\mathbf Q$} (rb);
\draw[sig] (rb.east) -- (11.9,0) node[sgn,right]{$\mathbf q,\dot{\mathbf q}$};
\draw[line width=0.6pt] (11.0,0) -- (11.0,-3.5);
\fill (11.0,0) circle (0.9pt); \fill (11.0,-1.3) circle (0.9pt); \fill (11.0,-2.5) circle (0.9pt); \fill (11.0,-3.5) circle (0.9pt);
\draw[sig] (11.0,-1.3) -- (be.east);
\draw[sig] (11.0,-2.5) -- (jf.east); \draw[sig] (jf.west) -- node[sgn,above,pos=0.2]{$\dot{\mathbf x}$} node[sgn,above,pos=0.93]{$-$} (sv.east);
\draw[sig] (11.0,-3.5) -- (kk.east); \draw[sig] (kk.west) -- node[sgn,above,pos=0.1]{$\mathbf x$} node[sgn,above,pos=0.96]{$-$} (sp.east);
\draw[sig] (i1.east) -- (1.05,-2.5) arc[start angle=180,end angle=0,radius=0.15] -- (sv.west); \node[sgn] at (2.35,-2.28) {$+$};
\draw[sig] (i0.east) -- (sp.west); \node[sgn] at (0.8,-3.28) {$+$};
\draw[sig] (sv.north) -- node[sgn,right]{$\dot{\mathbf E}$} (kv.south);
\draw[sig] (sp.north) -- node[sgn,left]{$\mathbf E$} (kp.south);
\draw[sig] (11.0,0) -- (11.0,1.0) -- (4.7,1.0) -- (mx.north); \fill (7.3,1.0) circle (0.9pt); \draw[sig] (7.3,1.0) -- (jt.north); \node[sgn,above] at (9.0,1.0) {$\mathbf q$};
\node[grpsv,fit=(ss)(kp)(kv)(sv)(sp)] (gsv) {};
\node[grplab,text=svcol,anchor=south west] at (gsv.north west) {SERVO};
\end{tikzpicture}}\par
\drawncap{Cartesian-based controller (slide p.\,33). Orange includes the kinematic model $\mathcal K,\mathbf J,\mathbf J^\top$; feedback $\mathbf q\to\mathcal K\to\mathbf x$, $\dot{\mathbf q}\to\mathbf J\to\dot{\mathbf x}$; $\mathbf q$ lines into $\mathbf M_x,\mathbf J^\top$ added.}
\end{wnfig}
\begin{worked}[$\mathbf Q=\mathbf J^\top\mathcal F$: planar 2R, $l_1=l_2=1$, $\mathbf q=(30^\circ,60^\circ)$]
$\mathbf J=\begin{bmatrix}-l_1s_1-l_2s_{12}&-l_2s_{12}\\l_1c_1+l_2c_{12}&l_2c_{12}\end{bmatrix}=\begin{bmatrix}-1.5&-1\\0.8660&0\end{bmatrix}$. $\mathcal F=(3,-4)$ N: $\mathbf Q=\mathbf J^\top\mathcal F=(-7.9641,\,-3.0)$ N\,m. Moments: tool $(0.866,1.5)$: $0.866(-4)-1.5(3)=-7.9641$; about joint 2 ($0,1$ away): $0(-4)-1(3)=-3$. \chk
\end{worked}
\begin{recipe}[Cartesian controller on paper]
\begin{enumerate}
\item Frame: $\mathbf x,\mathbf J,\mathcal F$ in $\{0\}$. \ 2.\ $\mathbf x=\mathcal K(\mathbf q)$, $\dot{\mathbf x}=\mathbf J\dot{\mathbf q}$; $\mathbf E_x,\dot{\mathbf E}_x$.
\item $\mathcal F=\mathbf M_x(\ddot{\mathbf x}_d+\mathbf K_v\dot{\mathbf E}_x+\mathbf K_p\mathbf E_x)+\mathbf c_x+\mathbf g_x$.
\item $\mathbf Q=\mathbf J^\top\mathcal F$ (transpose, $\mathbf J$ at the current $\mathbf q$).
\item Check: static $\Rightarrow\mathbf Q=\mathbf g$; torque $=\mathbf r\times\mathcal F$; singular poses (6-DoF: $\mathbf J$ is $6\times n$).
\end{enumerate}
\end{recipe}

\hd{Force control \pg{34--44}}
\premise{mass $m$ on a wall (stiffness $k_e$), unknown $f_{dist}$; $f_d(t)$; measured $f_e$, $\dot x$}{$f$ so $f_e\to f_d$}
\begin{deriv}[plant and partitioned law \pg{36--38}]\label{w9:force}
$f_e=k_ex\Rightarrow\ddot x=k_e^{-1}\ddot f_e$. Forces \pg{36}: $f-f_e-f_{dist}=m\ddot x$, so $f=mk_e^{-1}\ddot f_e+f_{dist}+f_e$. Controlled variable $f_e$: $\alpha=mk_e^{-1}$, $\beta=f_{dist}+f_e$, $E=f_d-f_e$:
\[f=mk_e^{-1}\big[\ddot f_d+k_v\dot E+k_pE\big]+f_{dist}+f_e\Rightarrow\ddot E+k_v\dot E+k_pE=0.\]
\textbf{Infeasible:} contains the unknown $f_{dist}$ \pg{38}.
\end{deriv}
\begin{fbx}[The two feasible options \pg{39--40}; $\alpha=mk_e^{-1}$, $X=\alpha k_p$]
\[\begin{gathered}\text{Opt.\,1 (drop }f_{dist}):\ f=\alpha[\ddot f_d+k_v\dot E+k_pE]+f_e;\\ \ddot E+k_v\dot E+k_pE=\frac{f_{dist}}{\alpha},\ \ E_{ss}=\frac{f_{dist}}{\alpha k_p}\end{gathered}\]
\[\begin{gathered}\text{Opt.\,2 (}f_d\text{ for }f_{dist}{+}f_e):\ f=\alpha[\ddot f_d+k_v\dot E+k_pE]+f_d;\\ \ddot E+k_v\dot E+\Big(\frac1\alpha+k_p\Big)E=\frac{f_{dist}}{\alpha},\ \ E_{ss}=\frac{f_{dist}}{1+\alpha k_p}\end{gathered}\]
$f_e=k_ex$ measured contact force; $\alpha$ in s$^2$. $E_2/E_1=X/(1+X)<1$: Opt.\,2 always better. Both non-zero unless $f_{dist}=0$. Opt.\,2 adds $k_e/m$ to the stiffness term.
\end{fbx}
\begin{deriv}[the two loops \added]
Plant $\alpha\ddot f_e+f_{dist}+f_e=f$. Opt.\,1: $\alpha(\ddot f_d-\ddot f_e)+\alpha k_v\dot E+\alpha k_pE=f_{dist}$. Opt.\,2: $\alpha\ddot E+\alpha k_v\dot E+\alpha k_pE=f_{dist}+f_e-f_d=f_{dist}-E\Rightarrow$ static $E=\dfrac{f_{dist}/\alpha}{k_p+1/\alpha}=\dfrac{f_{dist}}{1+\alpha k_p}$.
\end{deriv}
\begin{wnfig}
\resizebox{\linewidth}{!}{\begin{tikzpicture}[x=0.9cm,y=0.68cm]
\node[sgn,anchor=east] (i2) at (0,0) {$\ddot f_d$};
\node[sgn,anchor=east] (i1) at (0,-2.5) {$\dot f_d$};
\node[sgn,anchor=east] (i0) at (0,-3.5) {$f_d$};
\node[sumj] (ss) at (2.1,0) {};
\node[svblk,minimum width=9mm] (kp) at (1.2,-1.3) {$k_p$};
\node[svblk,minimum width=9mm] (kv) at (2.8,-1.3) {$k_v$};
\node[sumj] (sv) at (2.8,-2.5) {};
\node[sumj] (sp) at (1.2,-3.5) {};
\node[mbblk] (al) at (4.5,0) {$mk_e^{-1}$};
\node[sumj] (sb) at (5.9,0) {};
\node[plant,minimum width=31mm] (rb) at (8.7,0) {\footnotesize$f=m\ddot x{+}f_{dist}{+}k_ex$};
\node[blk,minimum width=10mm] (dd) at (8.7,-2.5) {$\tfrac{d}{dt}(\ )$};
\draw[sig] (i2.east) -- (ss.west); \node[sgn] at (1.6,0.22) {$+$};
\draw[sig] (kp.north) -- (ss.235); \node[sgn] at (1.35,-0.55) {$+$};
\draw[sig] (kv.north) -- (ss.305); \node[sgn] at (2.65,-0.55) {$+$};
\draw[sig] (ss) -- (al);
\draw[sig] (al) -- node[sgn,pos=0.45,above]{$+$} (sb);
\draw[sig] (sb) -- node[sgn,above]{$f$} (rb);
\draw[sig] (rb.east) -- (11.5,0) node[sgn,right]{$f_e$};
\draw[line width=0.6pt] (10.9,0) -- (10.9,-3.5); \fill (10.9,0) circle (0.9pt); \fill (10.9,-2.5) circle (0.9pt);
\draw[sig] (10.9,-2.5) -- (dd.east); \draw[sig] (dd.west) -- node[sgn,above,pos=0.3]{$\dot f_e$} node[sgn,above,pos=0.93]{$-$} (sv.east);
\draw[sig] (10.9,-3.5) -- node[sgn,above,pos=0.25]{$f_e$} node[sgn,above,pos=0.93]{$-$} (sp.east);
\draw[sig] (i1.east) -- (1.05,-2.5) arc[start angle=180,end angle=0,radius=0.15] -- (sv.west); \node[sgn] at (2.35,-2.28) {$+$};
\draw[sig] (i0.east) -- (sp.west); \node[sgn] at (0.8,-3.28) {$+$};
\fill (0.5,-3.5) circle (0.9pt);
\draw[sig] (0.5,-3.5) -- (0.5,-2.65) arc[start angle=-90,end angle=90,radius=0.15] -- (0.5,-0.15) arc[start angle=-90,end angle=90,radius=0.15] -- (0.5,1.0) -- (5.9,1.0) -- (sb.north); \node[sgn] at (6.12,0.45) {$+$}; \node[sgn,above] at (3.2,1.0) {$f_d$ (Option 2)};
\draw[sig] (sv.north) -- node[sgn,right]{$\dot E$} (kv.south);
\draw[sig] (sp.north) -- node[sgn,left]{$E$} (kp.south);
\node[grpsv,fit=(ss)(kp)(kv)(sv)(sp)] (gsv) {};
\node[grplab,text=svcol,anchor=south west] at (gsv.north west) {SERVO};
\end{tikzpicture}}\par
\drawncap{force controller (slide p.\,41). Top path adds $f_d$ = Option 2 (Option 1: measured $f_e$). The $\tfrac d{dt}$ block is flagged ``not a good idea'' (p.\,42): it amplifies force noise.}
\end{wnfig}
Practical law \pg{43--44}: constant $f_d$, $\dot E=-k_e\dot x$, so Opt.\,2 with $\dot f_e\to k_e\dot x$. $k_e$ is chosen high; design robust to small changes in $k_e$.
\begin{wnfig}
\resizebox{\linewidth}{!}{\begin{tikzpicture}[x=0.9cm,y=0.68cm]
\node[sgn,anchor=east] (fd) at (0,0) {$f_d$};
\node[sumj] (s1) at (1.4,0) {};
\node[mbblk,minimum width=11mm] (g1) at (3.1,0) {$\tfrac{k_p}{k_e}$};
\node[sumj] (s2) at (4.8,0) {};
\node[mbblk,minimum width=9mm] (m) at (6.0,0) {$m$};
\node[sumj] (s3) at (7.3,0) {};
\node[plant,minimum width=31mm] (rb) at (9.6,0) {\footnotesize$f=m\ddot x{+}f_{dist}{+}k_ex$};
\node[svblk,minimum width=9mm] (kv) at (4.8,-1.4) {$k_v$};
\draw[sig] (fd.east) -- (s1.west); \node[sgn] at (0.95,0.22) {$+$};
\draw[sig] (s1) -- node[sgn,above]{$E$} (g1);
\draw[sig] (g1) -- (s2); \node[sgn] at (4.3,0.22) {$+$};
\draw[sig] (s2) -- (m);
\draw[sig] (m) -- (s3); \node[sgn] at (6.85,0.22) {$+$};
\draw[sig] (s3) -- node[sgn,above]{$f$} (rb);
\draw[sig] (rb.east) -- (12.0,0) node[sgn,right]{$f_e$};
\fill (0.4,0) circle (0.9pt); \draw[sig] (0.4,0) -- (0.4,1.0) -- (7.3,1.0) -- (s3.north); \node[sgn] at (7.55,0.45) {$+$};
\draw[sig] (11.4,0) -- (11.4,-2.4) -- (1.4,-2.4) -- (s1.south); \fill (11.4,0) circle (0.9pt); \node[sgn] at (1.65,-0.5) {$-$}; \node[sgn,above] at (9,-2.4) {$f_e$};
\draw[sig] (10.3,-0.38) -- (10.3,-1.4) -- (kv.east); \node[sgn,right] at (10.3,-0.95) {$\dot x$};
\draw[sig] (kv.north) -- (s2.south); \node[sgn] at (5.05,-0.55) {$-$};
\end{tikzpicture}}\par
\drawncap{practical force controller (slide p.\,43): $f=m\big[k_e^{-1}k_pE-k_v\dot x\big]+f_d$; $\dot x$ measured, $f_e$ never differentiated.}
\end{wnfig}
\begin{worked}[$m=2$, $f_{dist}=5$, $k_p=100$, $k_v=20$]
$k_e=1000$: $\alpha=0.002$, $X=0.2$, $E_1=25$ N, $E_2=5/1.2=4.17$ N. Soft $k_e=20$: $X=10$, $E_1=0.5$, $E_2=0.455$. $E_2=1\,\%$ of $f_{dist}$ (stiff): $1+X=100$, $k_p=49\,500$. Practical law, $f_d=10$, $f_e=7$, $\dot x=0.003$: $E=3$, $f=2[0.001(100)(3)-20(0.003)]+10=10.48$ N $=$ Opt.\,2. Opt.\,2 loop $\omega_n^2=k_p+k_e/m=600$, critical $k_v=49.0$. \chk
\end{worked}
\begin{trap}
Opt.\,1 $E_{ss}$ used for Opt.\,2 (missing $1+$); ``stiffer wall helps'' (both grow with $k_e$; Opt.\,2 saturates at $f_{dist}$); differentiating $f_e$; $E=f_e-f_d$; $k_v=2\sqrt{k_p}$ for Opt.\,2.
\end{trap}

\hd{Hybrid force/position control \pg{45--47}}\label{w9:hyb}
\premise{some directions blocked (force), others free (position); $\ddot{\mathbf x}_d,\dot{\mathbf x}_d,\mathbf x_d$, $\mathcal F_d$; constraints}{$\mathcal F$, then $\mathbf Q=\mathbf J^\top\mathcal F$}
\begin{fbx}[Selection matrices \pg{47}]
\[\mathcal F=\mathbf S\,\mathcal F_{pos}+\mathbf S'\,\mathcal F_{force},\qquad \mathbf S'=\mathbf I-\mathbf S,\qquad \mathbf Q=\mathbf J^\top\mathcal F\]
$\mathbf S$ diagonal, entry \textbf{1 = position-controlled}; $\mathbf S'$ entry \textbf{1 = force-controlled}; $\mathcal F_{pos}$ from the position law (feedback $\mathbf x$), $\mathcal F_{force}$ from the force law ($\mathcal F_d$, measured $\mathcal F_E$). $\mathbf S'=\mathbf I-\mathbf S$ is my reading of ``complementary''. Write in the \emph{constraint frame} (rotate $\mathcal F$ if the surface is tilted) \added.
\end{fbx}
\begin{wnfig}
\resizebox{\linewidth}{!}{\begin{tikzpicture}[x=0.85cm,y=0.68cm]
\node[sgn,anchor=east] (a2) at (0,1.4) {$\ddot{\mathbf x}_d$};
\node[sgn,anchor=east] (a1) at (0,0.9) {$\dot{\mathbf x}_d$};
\node[sgn,anchor=east] (a0) at (0,0.4) {$\mathbf x_d$};
\node[svblk,minimum height=14mm,minimum width=18mm] (pl) at (2.2,0.9) {position\\control law};
\node[svblk,minimum height=12mm,minimum width=18mm] (fl) at (2.2,-1.8) {force\\control law};
\node[blk,minimum width=9mm] (S) at (5.0,0.9) {$\mathbf S$};
\node[blk,minimum width=9mm] (Sp) at (5.0,-1.8) {$\mathbf S'$};
\node[sumj] (sg) at (6.9,-0.45) {};
\node[plant,minimum height=14mm,minimum width=20mm] (rb) at (9.2,-0.45) {Cartesian\\arm};
\node[sgn,anchor=east] (cn) at (0.6,-0.45) {constraints};
\node[sgn,anchor=east] (fd) at (0,-1.8) {$\mathcal F_d$};
\draw[sig] (a2.east) -- (pl.west |- a2); \draw[sig] (a1.east) -- (pl.west |- a1); \draw[sig] (a0.east) -- (pl.west |- a0);
\draw[sig] (fd.east) -- (fl.west);
\draw[sig] (pl.east) -- (S.west); \draw[sig] (fl.east) -- (Sp.west);
\draw[line width=0.6pt] (cn.east) -- (5.0,-0.45); \fill (5.0,-0.45) circle (0.9pt);
\draw[sig] (5.0,-0.45) -- (S.south); \draw[sig] (5.0,-0.45) -- (Sp.north);
\draw[sig] (S.east) -- (sg.north west); \node[sgn] at (6.5,0.1) {$+$};
\draw[sig] (Sp.east) -- (sg.south west); \node[sgn] at (6.5,-1.0) {$+$};
\draw[sig] (sg) -- node[sgn,above]{$\mathcal F$} (rb);
\draw[sig] ($(rb.east)+(0,0.4)$) -- ++(0.8,0) node[sgn,above,pos=0.5]{$\mathbf x$} -- (11.1,2.3) -- (2.2,2.3) -- (pl.north);
\draw[sig] ($(rb.east)+(0,-0.4)$) -- ++(0.8,0) node[sgn,below,pos=0.5]{$\mathcal F_E$} -- (11.1,-3.1) -- (2.2,-3.1) -- (fl.south);
\node[sgn,above] at (5.0,1.35) {\scriptsize(1 = position)};
\node[sgn,below] at (5.0,-2.2) {\scriptsize(1 = force)};
\end{tikzpicture}}\par
\drawncap{hybrid controller (slide p.\,47, after Craig). Both laws are servo (blue); $\mathbf S,\mathbf S'$ pick per axis which law reaches the arm.}
\end{wnfig}
\textbf{Entries of $\mathbf S$:} blocked $\to$ force ($\mathbf S$ 0), free $\to$ position (1). Table, $z$ normal (p.\,46): $\operatorname{diag}(1,1,0)$; wall normal $x$: $\operatorname{diag}(0,1,1)$; drawer along $x$ \added: $\operatorname{diag}(1,0,0)$; peg in vertical hole $(x,y,z,\theta_x,\theta_y,\theta_z)$ \added: $\operatorname{diag}(0,0,1,0,0,1)$.
\begin{worked}[wipe a table, $z$ normal; $\mathcal F_{pos}=(12,-3,9)$, $\mathcal F_{force}=(1,2,10)$ N]
$\mathbf S=\operatorname{diag}(1,1,0)$, $\mathbf S'=\operatorname{diag}(0,0,1)$: $\mathcal F=(12,-3,0)+(0,0,10)=(12,-3,10)$ N; position $9$ and force $1,2$ discarded. Normal $x$: $\mathbf S=\operatorname{diag}(0,1,1)$ gives $(1,-3,9)$. Peg: position law $(1,2,-15,0.3,-0.2,0.8)$, force law $(-3,4,2,0.5,-0.6,0.1)$, $\mathbf S=\operatorname{diag}(0,0,1,0,0,1)$: $\mathcal F=(-3,4,-15,0.5,-0.6,0.8)$. \chk
\end{worked}
\begin{trap}
$\mathbf S,\mathbf S'$ swapped (eraser ignores the press); both 1 on one axis (laws fight); world axes on a tilted surface; forgetting $\mathbf Q=\mathbf J^\top\mathcal F$.
\end{trap}

\hd{Exam questions}
Only the practical sheet is lecturer-written (W13 Q4 EoM is the stepping stone: derive the EoM, then partition). Marks not published; indicative ~1 per step.
\begin{examq}[Practical sheet (a)--(c) \ \mk{marks assumed}]\label{w9:prac}
Arm masses $m=(3,2,1)$ kg at the distal link ends; rest at $(x,y,\varphi)=(5,1,90^\circ)$ to rest at $(-3,3,180^\circ)$. (a) Joint angles, cubics ($T=2$ s assumed). (b) Lagrange $\mathbf Q=\mathbf M\ddot{\mathbf q}+\mathbf C+\mathbf g$. (c) Evaluate at start.
\solution
\textbf{(a)} $\varphi=\theta_1+\theta_2+\theta_3$; wrist $\mathbf w=(x-L_3\cos\varphi,y-L_3\sin\varphi)$, then planar 2R. Start: $\mathbf w=(5,0)$, $\cos\theta_2=(25-9-4)/12=1$, $\theta_2=0$, $\theta_1=0$, $\theta_3=90^\circ$: $\mathbf q_0=(0,0,90^\circ)$; FK $(5,1)$ \checkmark. Goal: $\mathbf w=(-2,3)$, $|\mathbf w|^2=13$, $\cos\theta_2=0$, $\theta_2=90^\circ$; $\theta_1=\operatorname{atan2}(3,-2)-\operatorname{atan2}(2,3)=123.69^\circ-33.69^\circ=90^\circ$; $\theta_3=0$: $\mathbf q_f=(90^\circ,90^\circ,0)$. Other branch $(157.38^\circ,-90^\circ,112.62^\circ)$. \chk\\
Cubic $q=q_0+a_2t^2+a_3t^3$, $a_2=3\Delta/T^2$, $a_3=-2\Delta/T^3$, $\Delta=(\tfrac\pi2,\tfrac\pi2,-\tfrac\pi2)$: joints 1, 2: $a_2=1.1781$, $a_3=-0.3927$; joint 3: $-1.1781$, $+0.3927$. Peak speed $1.5\Delta/T=1.1781$ at $t=1$; peak acceleration $6\Delta/T^2=2.3562$ at $t=0$. Joint cubic is not a straight tool line: $t=1$ gives $(45^\circ,45^\circ,45^\circ)$, tool $(1.414,4.828)$. \chk
\end{examq}
\begin{examq}[Practical sheet (b), (c)]
Positions $\mathbf p_1=L_1(c_1,s_1)$, $\mathbf p_2=\mathbf p_1+L_2(c_{12},s_{12})$, $\mathbf p_3=\mathbf p_2+L_3(c_{123},s_{123})$; $K=\sum\tfrac12m_i|\dot{\mathbf p}_i|^2$, $P=g\sum m_iy_i$; $M_{ij}=\partial^2K/\partial\dot q_i\partial\dot q_j$, $g_i=\partial P/\partial q_i$. $A=18$, $B=3$, $D=2$ kg\,m$^2$:
\[\begin{aligned}
M_{11}&=67+36c_2+6c_{23}+4c_3, & M_{12}&=13+18c_2+3c_{23}+4c_3,\\
M_{13}&=1+3c_{23}+2c_3, & M_{22}&=13+4c_3,\\
M_{23}&=1+2c_3, & M_{33}&=1
\end{aligned}\]
\[\begin{aligned}
C_1&=-As_2(2\dot q_1\dot q_2+\dot q_2^2)-Bs_{23}\big(2\dot q_1(\dot q_2{+}\dot q_3)+(\dot q_2{+}\dot q_3)^2\big)\\&\quad-Ds_3\big(2(\dot q_1{+}\dot q_2)\dot q_3+\dot q_3^2\big)\\
C_2&=(As_2+Bs_{23})\dot q_1^2-Ds_3\big(2(\dot q_1{+}\dot q_2)\dot q_3+\dot q_3^2\big)\\ C_3&=Bs_{23}\dot q_1^2+Ds_3(\dot q_1+\dot q_2)^2
\end{aligned}\]
\[g_1=g[18c_1+6c_{12}+c_{123}],\ \ g_2=g[6c_{12}+c_{123}],\ \ g_3=g\,c_{123}\ \chk\]
\textbf{(c)} $\mathbf q_0=(0,0,90^\circ)$, $\dot{\mathbf q}=\mathbf 0$:
\[\mathbf M=\begin{bmatrix}103&31&1\\31&13&1\\1&1&1\end{bmatrix},\ \ \mathbf C=\mathbf 0,\ \ \mathbf g=9.81\begin{bmatrix}24\\6\\0\end{bmatrix}=\begin{bmatrix}235.44\\58.86\\0\end{bmatrix}\]
Check $g_1=9.81(3\cdot3+2\cdot5+1\cdot5)$ \checkmark; $\mathbf Q(0)=\mathbf M\ddot{\mathbf q}_d+\mathbf g=(548.8,160.2,2.4)$. \chk
\end{examq}
\begin{examq}[Practical sheet (d)--(i) \ \mk{marks assumed}]
(d)(e) Partitioned law, link 1. (g) Error along the path. (h) $\hat m=2.8$. (i) All joints.
\solution
\textbf{(e)} Link 1 alone: $\tau=m_1L_1^2\ddot q+m_1gL_1\cos q$; $\alpha=27$, $\beta=88.29\cos q$; $k_p=100$, $k_v=20$: $\tau=27(\ddot q_d+20\dot E+100E)+88.29\cos q$. At $t=0$: $\tau=27(2.3562)+88.29=151.91$ N\,m (joint 1 of the full arm: $\alpha=M_{11}=103$, $\beta=g_1=235.44$). \chk\\
\textbf{(g)} Exact model: $\ddot E+20\dot E+100E=0$, roots $-10,-10$; on the path $E\equiv0$; $0.1$ rad behind: $E(0.1)=0.0736$, no overshoot. Refresh $10$/$50$ ms: peak $0.26$/$1.13$ mrad. \chk\\
\textbf{(h)} $m=3$, $\hat m=2.8$, $L=3$: $\ddot E+20\dot E+100E=\tfrac1{14}[\ddot q+3.27\cos q]$. Horizontal: $E_{ss}=\tfrac1{14}(3.27)/100=2.34$ mrad $(0.134^\circ)$; vertical: $0$. On the cubic: forcing at $t=0$ $0.402$; simulated peak $3.35$ mrad at $t\approx0.47$ s. Full arm ($m_1=2.8$): $\hat M_{11}=101.2$, $\hat g_1=229.554$. \chk\\
\textbf{(i)} $\mathbf Q=\mathbf M(\ddot{\mathbf q}_d+20\dot{\mathbf E}+100\mathbf E)+\mathbf C+\mathbf g$, $\mathbf M,\mathbf C,\mathbf g$ from (b) each step: three decoupled critically damped loops. New vs one joint: $\boldsymbol\alpha$ full; recompute each step (outer loop). With $\operatorname{diag}(\mathbf M)$: row 1 forcing $(31(2.3562)-2.3562)/103=0.6863$ rad/s$^2$ (29\,\% of $2.36$). \chk
\end{examq}
\begin{examq}[EQ3: computed torque, 2R arm \ \mk{marks assumed: 8}]
2R, $l_1=l_2=1$, masses $1$ kg at link ends: $\mathbf M=\begin{bmatrix}3+2c_2&1+c_2\\1+c_2&1\end{bmatrix}$, $\mathbf C=\begin{bmatrix}-s_2(2\dot q_1\dot q_2+\dot q_2^2)\\s_2\dot q_1^2\end{bmatrix}$, $\mathbf g=g\begin{bmatrix}2c_1+c_{12}\\c_{12}\end{bmatrix}$. At $\mathbf q=(30^\circ,60^\circ)$, $\dot{\mathbf q}=(1,-0.5)$: $\ddot{\mathbf q}_d=(0.5,0)$, $\mathbf E=(0.02,-0.01)$, $\dot{\mathbf E}=\mathbf 0$, $\mathbf K_p=100\mathbf I$, $\mathbf K_v=20\mathbf I$. (a) Law [2]. (b) $\mathbf M,\mathbf C,\mathbf g$ [3]. (c) $\mathbf Q$ [2]. (d) $\acute{\mathbf Q}=(0.5,0)$: what does joint 2 still receive, why? [1]
\solution
(a) $\mathbf Q=\mathbf M\acute{\mathbf Q}+\mathbf C+\mathbf g$, $\acute{\mathbf Q}=\ddot{\mathbf q}_d+\mathbf K_v\dot{\mathbf E}+\mathbf K_p\mathbf E$.
(b) $c_2=0.5$, $s_2=0.866$, $c_1=0.866$, $c_{12}=0$: $\mathbf M=\begin{bmatrix}4&1.5\\1.5&1\end{bmatrix}$; $C_1=-0.866(2(1)(-0.5)+0.25)=0.6495$, $C_2=0.866$; $\mathbf g=9.81(1.732,0)=(16.991,0)$.
(c) $\acute{\mathbf Q}=(2.5,-1.0)$; $\mathbf M\acute{\mathbf Q}=(8.5,2.75)$; $\mathbf Q=\boxed{(26.14,\,3.62)\text{ N\,m}}$.
(d) $\mathbf M(0.5,0)^\top=(2.0,0.75)$: joint 2 gets $0.75$ N\,m, cancelling $M_{21}\ddot q_1$ ($\boldsymbol\alpha$ not diagonal). \chk
\end{examq}
\end{multicols}
